\subsection{Exact input definition and computation}\label{sup:input-algorithm}
\label{subsec:input}
\begin{definition}\label{def:input}
A presented field $k_0\subset\C$ has finite representations for its elements.
Arithmetic and equality are computable. Its embedding is specified.
Model (M-alg) supplies a number field $\Q(\theta)$.
A minimal polynomial and isolating disc specify $\theta$.
In model (M-gen), oracle (O1) takes any finite input tuple.
It returns a $\Z$-basis of its integer relations.
For number fields this oracle is supplied by integer linear algebra.
To test $a\in\Z$, ask whether there is a relation $ma+n=0$ with
$m\ne0$ and $m\mid n$.

The rates, nodes, scales and Kummer parameters belong to $k_0$.
Bessel orders are supplied as reduced rationals.
Rational kernels are supplied as polynomial pairs.
Model (M-gen) also supplies their linear denominator factors.
In (M-alg), we compute a splitting field.
Set $k_1=k_0(\text{poles})$. The value statements require (M-alg) and
$\alpha\in k_0^\times$. The supplied constants $\gamma_F$ encode
the chosen base-point normalization. Thus $z_*$ does not enter the computation.
A base-point primitive is specified by $\gamma_F\in k_0$:
\begin{equation}\label{eq:primitive-normalization}
F=\sum_{\sigma=(\mu,p)}e^{\mu p}\operatorname{res}(F)_\sigma\Phi_\sigma+r_F+\gamma_F.
\end{equation}
At nonzero weights, rational reduction uniquely determines $r_F$.
At weight zero, its polynomial part must have zero constant.
Thus $\gamma_L=0$. Coordinates of type (B2$^E$) and towers instead use
origin normalization. Let $\mathcal V_{{\rm res},0}$ be the origin-normalized first-level residue span.
Write $\mathcal V_{{\rm res},*}$ for the base-point-normalized first-level residue span.
The list is normalization-compatible when
\begin{equation}\label{eq:compatibility}
\mathcal V_{{\rm res},0}\cap\mathcal V_{{\rm res},*}=0,
\end{equation}
or after moving a specified origin subset to the base-point side using
supplied $\gamma_G\in k_0$ satisfying \eqref{eq:primitive-normalization}.
A supplied comparison is assumed correct; its constant is exact input.
The algorithm does not evaluate the analytic primitive at the base point.
Choose residue bases separately on the two
normalization sides. Each fundamental coordinate-basis relation uses one normalization.
An arbitrary relation decomposes into the two types.
Compatibility holds if all coordinates except possibly $L$ are origin-normalized.
Their residue vectors have $\sum_\mu v_{(\mu,0)}=0$.
For $L$, $\operatorname{res}(L)=e_{(0,0)}$.
It also holds if every coordinate uses the base point.
\end{definition}

\begin{proposition}\label{prop:algorithm}
In both input models, the quantities $d,\rho,\tau,\#\mathcal C$ and the
corresponding degrees are computable. For compatible lists, the generators
and every sublist ideal are computable over $k_1(z)[\Theta^{\pm1}]$.
Model (M-alg) requires no oracle, and (M-gen) uses (O1). For incompatible
lists, the algorithm reports undetermined mixed constants and returns no
ideal. Inputs supplied only by complex approximations are excluded.
For values, the hypotheses of main-text Theorem~\ref{thm:values} are required.
\end{proposition}
\begin{proof}

The algorithm first rejects inadmissible specifications. It then closes
and normalizes the list, retaining one representative per coordinate.
It emits $X_q-X_{\operatorname{rep}(q)}$ for every duplicate slot of every type.
It emits $X_c-c$ for each listed constant $c\in k_1$.
Only distinct higher levels are free.
Half-integer Bessel coordinates retain main-text Section~\ref{sec:main-statement}'s exponential graphs.
They also retain its derivative/order-shift graphs.
The saturated lattice contains their rates $\pm c$.
They contribute no class to $\mathcal C$.
The algorithm preserves these graphs through closure and specialization.
It also preserves them through saturation and contraction.
It composes the linkage maps. Oracle (O1), followed by
Hermite normal form, gives the minimal saturated lattice. Partial
fractions and finite Taylor expansions determine the residues. Linear algebra then gives $\rho$, $\kappa_{ji}$ and the test for
$e_{(0,0)}\in\mathcal V_{\rm res}$; $\tau=\sum_\eta(h_\eta-1)$ is the number of distinct retained levels $(\eta,k)$ with $k\ge2$. Rational reduction with the specified
normalization gives
\begin{equation}\label{eq:relation-constants}
c_j=\begin{cases}\gamma_j-\sum_i\kappa_{ji}\gamma_i,&\text{base-point},\\
-\lim_0r_j(z,e^{\lambda_1z},\ldots,e^{\lambda_dz}),&\text{origin}.
\end{cases}
\end{equation}
Any supplied normalization comparisons are applied before this step.
Next adjoin inverse variables with relations
$\Theta_h\Theta_h^{-1}-1$, saturate, and eliminate. Contraction gives
the sublist ideals and the degrees of the original coordinate lists.
For values, specialize only $z=\alpha$.
The procedure uses field arithmetic and equality, with (O1).
It also uses Hermite normal form and Buchberger elimination.
For compatible lists, these operations obtain all required constants.
No analytic evaluation is used.
\end{proof}

The arithmetic coefficient bounds are printed in main-text Appendix~\ref{app:arithmetic}.
